\begin{frame}
  \frametitle{Formalization in Rocq using Autosubst}

  Why did we use the Autosubst 1?
  \begin{itemize}
  \item<2-> A library instead of a code generator.
  \item<3-> Simple to use and learn (built using $\sigma$-calculus).
  \item<4-> Minimalistic architecture which allows customizability.
  \end{itemize}
  
  What is missing (for us)?
  \begin{itemize}
  \item<5-> More advanced tactics to derive substituion operations (how?).
  \item<6-> Modularization for the substitution lemmas.
  \item<7-> Support for interactions between multiple calculi.
  % \item<8-> Of course support for mutually dependent syntactic classes would be good...
  \end{itemize}
  
\end{frame}

%%% Local Variables:
%%% mode: LaTeX
%%% TeX-master: "../presentation"
%%% End:
